\chapter{知识复用的判断与选择}
\label{chap:ket-reuse-selection}

花卉识别系统要进入一个新地区。手中已有一个物种分类模型、一个病害识别模型和一个花朵定位模型，目标端只能标注少量照片。该选来源任务准确率最高的模型，还是照片风格最接近的模型？若某个模型很快达到可用精度，却在充分训练后不如从头训练，它究竟有没有帮助？若两个单独表现普通的来源恰好补上彼此的不足，只按单项排名选模型又会漏掉什么？

这些问题都在问已有经验的价值，但需要三种不同的判断。收益判断规定怎样才算有帮助；可复用知识的判断考察已有信息与目标要求是否兼容；选择则在证据、访问条件和预算约束下决定采用哪些内容。第~\ref{chap:ket-overview}~章已经区分知识载体与任务能力，本章沿用这一点：被选择的是具体样本、表示、参数、模块、预测行为或学习规则，其价值取决于目标和使用方式，不能由来源模型的名字永久确定。

下文以来源联合分布$\mathcal P_S$、目标联合分布$\mathcal P_T$和目标损失$\ell_T$描述问题。目标标签既可能用于学习，也可能用于挑选来源，必须记录其用途。本章所有小规模数值、表格中的假设成绩及花卉任务关系均为\term{教学构造}，用于复算决策，不是论文实验或真实系统测量；引用的论文提供方法和证据边界。

\section{知识复用的收益判断}
\label{sec:ket-reuse-benefits}

“复用更好”只有在比较对象明确后才有意义。若复用方案获得更多目标标签、更大的模型和更多调参机会，最终的差异就不能全部归给来源知识。这里把无来源起点的学习方案称为\term{无迁移对照}（no-transfer baseline），即在约定的信息条件下，不使用待检验来源知识而完成同一目标任务的方案。对照仍应使用合理的训练程序，不能故意缩短训练或放弃必要调参来制造收益。

固定比较条件并不要求把所有方案强行做成同一种训练过程。冻结表示与全量微调本来就改变不同数量的参数，但应交代这种差异。若问题是来源初始化本身是否有效，尽量只改变初始化，保持目标模型结构、数据和优化搜索范围一致；若问题是在设备预算内采用哪种完整方案，则应比较各自可执行的最佳配置，同时把结论称为方案收益。

\subsection{预测效果收益的判断}
\label{sec:ket-reuse-prediction}

令学习程序$\mathcal A$接收来源知识$K$和目标训练集$S_T$，得到预测器$f_K=\mathcal A(K,S_T)$。关闭来源通道后得到$f_0=\mathcal A(\varnothing,S_T)$。目标风险与收益写为
\begin{equation}
  \begin{aligned}
    R_T(f)&=\mathbb E_{(\bm{x},y)\sim\mathcal P_T}
      [\ell_T(f(\bm{x}),y)],\\
    \Delta_T(K)&=R_T(f_0)-R_T(f_K).
  \end{aligned}
  \label{eq:ket-reuse-risk-gain}
\end{equation}
在损失越小越好的约定下，$\Delta_T(K)>0$表示\term{正迁移}（positive transfer），即来源使目标风险下降；$\Delta_T(K)<0$表示\term{负迁移}（negative transfer），即使用来源反而增加目标风险。若学习过程有随机性，应对训练集抽样和初始化等随机因素取期望，实际用重复运行估计，不能把一次波动直接解释为稳定的迁移方向。

Wang等将负迁移明确为同一学习算法有无来源信息时的目标风险差，强调算法、来源与目标的关系以及目标标签量都影响结论。式~\eqref{eq:ket-reuse-risk-gain}~用收益而不是风险增量记号，因此符号与该文的负迁移间隙相反。这个定义也说明，一个迁移算法输给另一个更强的无迁移算法，值得在实践中警惕，却还不能单凭这一比较把全部差距归因于来源知识\scite{ket-wang2019-negative-transfer}

风险一般不可直接获得。对未参与拟合的目标样本集$V_T$，可计算经验风险$\hat R_{V_T}(f)=|V_T|^{-1}\sum_{(\bm{x},y)\in V_T}\ell_T(f(\bm{x}),y)$。二分类的零一损失对应错误率，准确率收益便等于两个错误率之差；病害漏检成本高时，应改用相应损失或同时报告召回率，不能用整体准确率替代既定目标。差异若只有一两个样本，还必须考虑估计不确定性。

表~\ref{tab:ket-reuse-checkpoints}~固定一个两物种目标任务、一组含100张照片的验证集、相同的目标网络结构和目标训练数据。冻结表示只训练新预测头；普通微调从来源参数开始更新全网络；从头训练使用随机初始化。三者使用相同的超参数搜索次数，表中列出一次假设运行在两个预先指定检查点的错误数。训练步数相同并不等于计算量相同，此处仅比较训练进度上的效果，资源另行核算。

\begin{table}[htbp]
  \centering
  \small
  \begin{tabularx}{\linewidth}{Xrr}
    \toprule
    方案 & 第10步错误数 & 第100步错误数 \\
    \midrule
    冻结来源表示，训练新头 & 18 & 16 \\
    来源初始化，全量微调 & 12 & 14 \\
    随机初始化，从头训练 & 30 & 10 \\
    \bottomrule
  \end{tabularx}
  \caption{同一目标任务上的教学构造。各列分母均为100张验证照片，检查点事先固定；这些数值不表示任何真实模型的学习曲线。}
  \label{tab:ket-reuse-checkpoints}
\end{table}

在第10步，普通微调相对从头训练的错误率收益为$(30-12)/100=0.18$，即18个百分点；在第100步则为$(10-14)/100=-0.04$，即负4个百分点。早期优势并没有保证训练终点优势。若上线期限只允许10步，前一种收益可能正是所需；若允许100步，依据早期成绩锁定方案就可能作出相反选择。微调曲线后段变差也可能来自过拟合或学习率不合适，表中数值只展示这种可能性，不证明某个具体原因。

冻结方案回答“来源表示中的信息现在能否被新头利用”，微调回答“允许调整表示后能否形成目标能力”，从头训练回答“没有该来源时能学到什么”。三种答案不应合并成一个模型通用性分数。比较冻结来源表示与冻结随机表示，可以进一步隔离来源训练对特征的影响；比较相同结构的来源初始化与随机初始化，则更接近初始化收益。置信区间、配对重复运行及训练终点的选择规则由第~\ref{chap:ket-evaluation}~章统一说明。

\subsection{样本效率收益的判断}
\label{sec:ket-reuse-sample-efficiency}

目标效果达到可用程度以后，还可能关心用了多少新证据。对于标签预算$n$，记在同一采样协议下训练的模型为$f_{K,n}$。把$n$与验证效果对应起来得到\term{学习曲线}（learning curve），即观察增加目标样本怎样改变效果。给定可接受风险$r_\star$，在实际考察的预算集合$N$上定义阈值到达量
\begin{equation}
  n_\star(K;r_\star)
  =
  \begin{cases}
    \min\{n\in N:\hat R_{V_T}(f_{K,n})\leq r_\star\},
      & \text{若集合非空},\\
    +\infty, & \text{若集合为空}.
  \end{cases}
  \label{eq:ket-reuse-label-threshold}
\end{equation}
这里的$+\infty$只表示在已测预算集合内没有到达阈值；报告时还应给出最大已测预算，而不能解释为任意预算下都无法到达。这个量依赖数据划分、类别构成、训练程序和阈值。真实风险未知时，它是经验到达量，不是达到总体风险要求的证明。学习曲线有随机波动时，还应事先说明是否要求多次运行的平均值、保守区间或连续几个预算点都满足阈值。

图~\ref{fig:ket-reuse-learning-curves}~给出可复算的教学构造：训练标签数依次为20、40、80、160，迁移方案在同一100张验证照片上正确预测82、88、91、92张，从头训练正确预测65、76、85、91张。所有类别在每个预算点都有样本，较小训练集嵌套于较大训练集。以准确率90\%为阈值，已测预算中的到达量分别为80与160，相差80个训练标签；在80个标签处比较，则是91\%与85\%的6个百分点效果差。前者固定效果比较标签，后者固定标签比较效果，不能把两种数字当成同一种收益。

\begin{figure}[htbp]
  \centering
  % Teaching construction; no measured results.
% Training labels: 20, 40, 80, 160.
% Correct predictions out of 100: reuse 82, 88, 91, 92;
% scratch 65, 76, 85, 91. No interpolation claims or uncertainty estimates.
% Coordinates: x = labels * 0.045 cm; y = (accuracy - 60) * 0.12 cm.
\begin{tikzpicture}[
  x=1cm,y=1cm,font=\figurefont\footnotesize,
  line width=0.8pt,
  axis/.style={draw=BookInk,-{Latex[length=1.8mm]}},
  grid/.style={draw=BookRule,line width=0.5pt},
  reuse/.style={draw=BookTeal,line width=1.1pt},
  scratch/.style={draw=BookBlue,dashed,line width=1.1pt}
]
  \foreach \v/\yy in {60/0,70/1.2,80/2.4,90/3.6,100/4.8} {
    \draw[grid] (0,\yy) -- (7.55,\yy);
    \node[anchor=east,text=BookInk] at (-0.12,\yy) {\v};
  }
  \draw[BookAmber,densely dotted,line width=0.9pt]
    (0,3.6) -- (7.55,3.6);
  \draw[axis] (0,0) -- (7.85,0);
  \draw[axis] (0,0) -- (0,5.05);
  \foreach \n/\xx in {20/0.9,40/1.8,80/3.6,160/7.2} {
    \draw[BookInk] (\xx,0) -- (\xx,-0.08);
    \node[anchor=north,text=BookInk] at (\xx,-0.12) {\n};
  }
  \node[text=BookInk] at (3.9,-0.7) {目标训练标签数（个）};
  \node[rotate=90,text=BookInk] at (-0.85,2.4) {准确率（\%）};
  \draw[reuse] plot coordinates
    {(0.9,2.64) (1.8,3.36) (3.6,3.72) (7.2,3.84)};
  \foreach \xx/\yy in {0.9/2.64,1.8/3.36,3.6/3.72,7.2/3.84} {
    \fill[BookTeal] (\xx,\yy) circle[radius=2.2pt];
  }
  \draw[scratch] plot coordinates
    {(0.9,0.6) (1.8,1.92) (3.6,3.0) (7.2,3.72)};
  \foreach \xx/\yy in {0.9/0.6,1.8/1.92,3.6/3.0,7.2/3.72} {
    \draw[BookBlue,fill=white,line width=0.8pt]
      (\xx-0.06,\yy-0.06) rectangle (\xx+0.06,\yy+0.06);
  }
  \draw[reuse] (0.6,5.4) -- (1.25,5.4);
  \fill[BookTeal] (0.925,5.4) circle[radius=2.2pt];
  \node[anchor=west,text=BookInk] at (1.35,5.4) {复用};
  \draw[scratch] (3.0,5.4) -- (3.65,5.4);
  \draw[BookBlue,fill=white] (3.265,5.34) rectangle (3.385,5.46);
  \node[anchor=west,text=BookInk] at (3.75,5.4) {从头训练};
  \node[anchor=west,text=BookAmber] at (5.5,5.4) {阈值90\%};
\end{tikzpicture}

  \caption{目标训练标签数与准确率的教学构造。点是明确给定的预算和正确数，连线仅引导阅读，不代表测得中间预算的结果。90\%虚线对应阈值，已测到达量为80和160；验证标签不计在横轴中，但须进入总标注账本。}
  \label{fig:ket-reuse-learning-curves}
\end{figure}

若某方案到160个标签仍未过线，应写“在已测预算内未达到”，或在记录中使用$>160$，而不是把160填成到达量。两个点之间也不能未经测量就宣称精确阈值位置。某个预算没有病害阳性样本，另一个预算覆盖了全部病害类别时，仅比较总条数会掩盖监督信息的差异；应补充每类数量、标注粒度和覆盖情况。一张照片的物种标签、一组像素掩码和一次人工交互反馈的成本不相等。

\begin{samepage}
标签账本还需要包括选择来源的消耗。假设迁移方案除80个适应标签外，还用了30个不重叠的来源筛选标签和100个验证标签，其目标开发标签总数为210；无迁移方案用160个训练标签和相同100个验证标签，总数为260。此时节省的是50个开发标签，不是只看适应阶段得到的80个。若同一批30个标签随后也用于适应，应按唯一标注对象计费，不重复加总，但记录它被用于两种决策，不能把复用后的训练效果当作独立验证。最终测试标签另外列示。

\end{samepage}

人工标签、无标签照片和交互轮次宜分别记账，组成预算向量
\[
  (n_{\mathrm{lab}},n_{\mathrm{unlab}},n_{\mathrm{feedback}}).
\]
三个分量依次表示人工标签数、无标签照片数和交互轮次；除非已明确单位成本，不能直接相加。主动查询挑出的困难样本和均匀抽样得到的支持集可能在相同条数下提供不同信息，查询策略及未被选中的样本池也属于协议。对少样本任务，还需固定支持集的类别数和每类样本数，否则“每次只用五张”可能描述完全不同的学习问题。

如果没有目标标签，只观察到输入分布，就无法一般地辨认来源是否降低了目标风险。来源与目标可能有完全相同的照片分布，却采用相反的标签规则；两种情形给无标签评分程序的输入一样，需要的决策却不同。第~\ref{chap:transfer-learning-theory}~章的不可识别性讨论说明了这种边界。无标签一致性或分布距离可以作为假设下的筛选依据，不能替代目标收益证据。

\subsection{计算与存储收益的判断}
\label{sec:ket-reuse-resource-benefits}

少用目标标签未必少用计算。获取来源模型、筛选候选、提取特征和搜索超参数都要消耗资源，部署时还可能保留多个模型。为避免把这些成本混在一起，令$C_{\mathrm{prep}}$为一次性来源准备成本，$C_{\mathrm{sel},k}$和$C_{\mathrm{adapt},k}$分别为目标任务$k$的来源选择与适应成本，$q_kc_{\mathrm{infer},k}$为该任务服务$q_k$次请求的推理成本。总计算与按任务摊销的准备成本为
\begin{equation}
  \begin{aligned}
    C_{\mathrm{reuse}}(m)
    &=C_{\mathrm{prep}}+
      \sum_{k=1}^{m}
      \bigl(C_{\mathrm{sel},k}+C_{\mathrm{adapt},k}
      +q_kc_{\mathrm{infer},k}\bigr),\\
    \overline C_{\mathrm{prep}}(m)&=\frac{C_{\mathrm{prep}}}{m}.
  \end{aligned}
  \label{eq:ket-reuse-amortization}
\end{equation}
这里$m$是实际使用该来源方案的目标任务数，不是为使平均值好看而设想的无限未来任务数。失败的迁移也消耗选择与适应成本，不能因其没有收益便从账本中删去；是否获益须另与无迁移对照比较。不同硬件上的GPU小时也不能直接当成等价计算量，应固定设备和计费条件，或分别报告设备时长、运算量与金额。

考虑一个教学构造：准备来源模型用120 GPU小时；每个目标选择来源用2 GPU小时，适应用8 GPU小时；每个任务从头训练用30 GPU小时。两条路线达到相同的目标效果，推理成本暂时相同，因而在比较中相消。复用与从头训练的总成本分别为$120+10m$和$30m$。只有
\begin{equation}
  120+10m\leq30m
  \quad\Longleftrightarrow\quad m\geq6
  \label{eq:ket-reuse-break-even}
\end{equation}
时，来源投入才能被目标端节省抵消。只服务一次任务，复用花130而非30 GPU小时；服务6次，两者都是180；服务10次，复用为220、从头训练为300，节省80 GPU小时。若复用服务每个任务还多花5 GPU小时推理，则盈亏平衡点改为$m\geq8$。准备投入相同，服务期限的选择仍会改变答案。

若一个已经训练好的模型可以依法直接取得，从当前项目的增量决策看，历史120 GPU小时可能是沉没成本，不应再次作为当前支出。此时仍要计入下载、格式转换、验证和授权等实际新增成本；从整个模型生命周期看，来源训练又不能凭空消失。报告应分别写清“已有来源可用的增量成本”与“包含来源准备的总成本”，不能只给有利的一种口径。

可训练参数少，也不等于存储和推理便宜。设骨干有1亿个参数，每个目标头有100万个参数，均用每参数4字节保存。忽略文件元数据时，十个全量微调模型需要$10\times404=4040$ MB；十个任务共用冻结骨干、各保存一个目标头，需要$400+10\times4=440$ MB。这里采用十进制MB，且不另存归档模型。共享节省了权重存储，但每次新照片仍须运行骨干。若为$n$张照片缓存$d$维特征，每元素$b$字节，缓存另占$ndb$字节；训练时的梯度、优化器状态和激活内存也应另计。保留多个来源共同推理时，这些负担可能重新增加。

人工核验、标签制作、计算和存储应分别列示。只有给出设备单价、存储保留期、请求量和人工计价依据，才可以合并成货币成本或效用；把准确率、GPU小时和MB直接相加没有解释意义。表~\ref{tab:ket-reuse-benefit-protocol}~把同一花卉系统的三种收益口径放在一起：一种预算下的优势，不自动延伸到另一种预算或另一个使用期限。

\begin{table*}[htbp]
  \centering
  \small
  \begin{tabularx}{\linewidth}{lXXXXX}
    \toprule
    收益 & 固定条件 & 比较量 & 必要对照 & 准备成本 & 常见误判 \\
    \midrule
    预测效果 & 目标标签、结构及调参机会 & 同一终点的错误率差 & 来源初始化与随机初始化；冻结随机特征可作补充 & 来源训练及目标端选择分别记录 & 第10步领先18个百分点，误称第100步仍领先 \\
    样本效率 & 效果阈值、类别覆盖与标注粒度 & 达阈值的标签、无标签及反馈预算 & 从头训练的完整学习曲线 & 筛选30个标签不能漏记 & 把80个适应标签当作全部标注支出 \\
    计算存储 & 达成效果、设备、请求量和期限 & 总计算、峰值内存、权重与缓存 & 同效果的从头训练或其他可执行方案 & 来源120 GPU小时按实际任务数摊销 & 一次任务更贵，却借十次任务的平均值声称省算力 \\
    \bottomrule
  \end{tabularx}
  \caption{知识复用收益的比较口径。数字来自本章的教学构造；各行回答不同问题，不把不同效果或不同成本边界的成绩混为一项排名。}
  \label{tab:ket-reuse-benefit-protocol}
\end{table*}

\section{可复用知识的判断}
\label{sec:ket-reuse-compatibility}

有了收益标准，仍不能直接从候选名单中选出赢家。一个来源可能覆盖目标照片，却没有保留病害纹理；也可能保留了纹理，但只允许查询物种类别，无法取得有用的中间表示。以下从输入输出关系、表示与预测关系、任务结构关系检查同一候选。这三种依据可以互相补充，不能把“分布接近”“标签含义相近”和“适合共用预测规则”当作同一句话。

\subsection{输入输出关系的判断}
\label{sec:ket-reuse-input-output}

输入兼容至少有两层含义。文件能读入模型，只说明尺寸、通道和编码可接受；观测兼容还要求这些数值表达相应的物理或语义对象。RGB照片与近红外图像都可能保存成三通道数组，却不能仅凭形状一致认定来源卷积核的响应仍有相同含义。来源拍摄地区只出现花朵正面，目标还包含严重遮挡或夜间照片时，需要指出哪些目标区域缺少来源覆盖，而不是用一个全局平均距离掩盖缺口。

输入分布相同也不足以保证输出规则相同。构造四种等概率照片属性，令$\bm{x}=(x_1,x_2)\in\{-1,1\}^2$，其中$x_1$表示花瓣颜色属性，$x_2$表示斑点属性。来源与目标都均匀采样这四种组合。来源标签为$y_S=\mathbb{I}(x_1=1)$，目标标签为$y_T=\mathbb{I}(x_2=1)$，其中$\mathbb{I}$为条件成立取1、否则取0的指示函数。直接复制来源预测规则，在目标上恰有两种组合答错，错误率为$1/2$；若目标改为$y_T=1-y_S$，直接复制则全部答错。两种情况下输入分布距离都为零，问题出在条件标签关系。

上述构造不表示来源表示一定无用。保留$(x_1,x_2)$的表示仍含有目标所需信息；只保留$x_1$的表示则无法判断独立的$x_2$。判断要具体到复用对象：源预测规则不适用，与源中间层完全没有价值，是两个不同结论。标注标准变化同样如此，例如来源把疑似病斑归为健康，目标规定必须送检；即使标签仍写“健康／异常”，决策边界也已经改变。

还要区分不兼容与缺少证据。假设来源只见过物种甲、乙，目标新增物种丙。来源中没有丙类照片，意味着对丙的覆盖证据缺失，不能直接推出来源边缘、纹理或形态表示有害；来源若把已确认的丙按旧标准标成乙，则出现了需要处理的标签冲突。这两种情况对应不同后续工作：前者需要补充目标证据或未知类处理，后者需要确认类别映射与标注规范。

可把检查结果落为三类可检验记录：目标哪些区域在来源中没有观察，哪些标签映射已由规范或配对样本确认，哪些观测对应仍没有依据。例如“目标夜间样本尚无来源覆盖；甲、乙的物种定义保持不变；丙不在来源输出集合内”比“两个领域很相似”更能指导选择。图~\ref{fig:ket-reuse-coverage-relations}~左侧只表达这种覆盖关系，不把区域重叠解释成标签机制一致。

覆盖和共同可解性的完整理论条件见第~\ref{chap:transfer-learning-theory}~章，尤其第~\ref{sec:adaptation-coverage}~节。需要怎样校正分布、建立跨观测对应，由第~\ref{chap:ket-update-maintenance}~章展开；新类别与新目标的语义要求由第~\ref{chap:ket-transfer-adaptation}~章展开。这里的产出是适用条件和证据缺口，还不是已经完成的适应。

\subsection{表示与预测关系的判断}
\label{sec:ket-reuse-representation}

来源网络的不同层可能保留不同信息。为了判断某一层能否支持目标任务，可以先固定特征提取器$\phi_\ell$，只训练一个容量受控的预测头。这称为\term{表示探测}（representation probing），即用约定难度的读取器检验特征中哪些信息能够被取出。它测量的是“这种读取方式下可用的信息”，不是对表示所含全部信息的穷尽判断。

设$\bm z_i=\phi_\ell(\bm x_i)\in\mathbb R^{d_\ell}$，目标有$c$类，一个线性探测头可由下式拟合：
\begin{equation}
  \begin{aligned}
    (\hat{\bm W},\hat{\bm b})
    =\argmin_{\bm W,\bm b}\Bigl\{
    &-\frac{1}{n}\sum_{i=1}^{n}
    \log\bigl[\operatorname{softmax}
    (\bm W\bm z_i+\bm b)\bigr]_{y_i}\\
    &+\lambda\lVert\bm W\rVert_F^2\Bigr\},
  \end{aligned}
  \label{eq:ket-reuse-linear-probe}
\end{equation}
其中$\bm W\in\mathbb R^{c\times d_\ell}$、$\bm b\in\mathbb R^c$，$\lVert\cdot\rVert_F$为矩阵Frobenius范数。来源参数不更新，正则强度$\lambda$及优化停止规则由独立验证过程确定。新照片先经同一预处理和$\phi_\ell$，再由新头输出目标概率；若缓存特征，提取器必须处于固定状态，不能在不同候选间不一致地更新归一化统计量。

探测成绩低有多种解释：表示丢失了目标信息，信息仍在但非线性读取才可取得，标签太少，或预测头没有优化好。增加探测器容量会改变问题；一个很大的读取器可能自己重新学习目标结构，弱化来源表示的贡献。比较不同层时，应记录特征维数、池化方式、头的参数量、标签量和优化预算，并用相同目标划分验证，而不是只写“训练一个分类器”。

回到四种等概率属性的教学构造。表示$\phi_A(\bm x)=(x_1,x_2)$与$\phi_B(\bm x)=x_1$都能把来源标签$y_S=\mathbb I(x_1=1)$完全分开，来源准确率同为100\%。目标若预测$y_T=\mathbb I(x_2=1)$，线性头可以从$\phi_A$完全读出标签；$\phi_B$把不同病害状态映成同一个值，对每个值的两种目标标签各占一半，任何只看$\phi_B$的预测器最高准确率都是50\%。这里不是优化失败，而是信息确实被合并了。来源成绩相同，没有提供区分这两种目标可用性的依据。

Yosinski等通过迁移不同层并比较冻结、微调与同任务重新连接的对照，区分了特征对来源任务的特异性和切断相邻层协同关系带来的优化困难。这一结果提示，复制一个模块后性能下降，不必只归因于模块本身不含有用信息；下游网络是否能够重新利用它也需要对照检验。该研究的证据来自所研究的卷积网络与任务划分，不能据此规定所有网络的某一固定层必然最通用\scite{yosinski2014}

访问接口会进一步限制判断。只有黑盒预测时，可以检查类别含义、概率是否归一化、置信度与目标标签的关系，却不能假装已观察到中间特征的线性可分性。来源头丢弃了病害信息，不意味着更早的表示也丢弃了；只是当前接口未提供验证途径。反过来，取得特征也不代表获得修改源参数的权限。第~\ref{sec:ket-reuse-fixed-score}~节的两种评分方法正好对应这两类接口，实际负迁移仍需用目标对照确认。

\subsection{任务结构关系的判断}
\label{sec:ket-reuse-task-relations}

来源和目标还可能在更高层共享结构。物种识别需要区分花朵形态，病害判断需要定位局部异常，花朵定位要求输出空间位置；三者输入同为照片，输出却分别是类别、病害状态和位置。定位提供的前景信息可能减少物种分类对背景的依赖，即使定位本身不是最终交付任务，它也可能作为辅助来源有用。是否有帮助，要看这种结构能否通过当前表示和学习程序被利用。

为把关系写成可测对象，令$f_{a\to b}$为从任务$a$的知识出发、按规定预算适应任务$b$的模型，$f_{0,b}$为任务$b$的无迁移对照，定义
\begin{equation}
  G_{a\to b}=R_b(f_{0,b})-R_b(f_{a\to b}).
  \label{eq:ket-reuse-directed-gain}
\end{equation}
该关系有方向，通常没有$G_{a\to b}=G_{b\to a}$的理由。两边不仅使用不同目标损失和对照，来源学习时保留的信息也不同。若两个任务指标的单位不一致，不能直接把$G$相加或用箭头粗细比较大小。

图~\ref{fig:ket-reuse-coverage-relations}~右侧给出一个固定预算下的教学构造。物种与病害各在100张目标验证照片上计错：病害从头训练错20张，物种来源适应后错18张，因此物种$\to$病害的收益为2个百分点；物种从头训练错10张，病害来源适应后错11张，反向收益为负1个百分点。定位来源适应物种后错6张，故定位$\to$物种为4个百分点。其余方向未测，图中不画箭头，绝不把缺测当成零收益。

\begin{figure*}[htbp]
  \centering
  % Teaching construction, not a measured distribution or task graph.
% Coverage rectangles encode set membership only, not probability mass.
% Directed gains (percentage points):
% species -> disease: (20 - 18) / 100 = +2;
% disease -> species: (10 - 11) / 100 = -1;
% localization -> species: (10 - 6) / 100 = +4.
% All other directions are unmeasured, not zero.
\begin{tikzpicture}[
  x=1cm,y=1cm,font=\figurefont\small,
  task/.style={
    rectangle,draw=BookInk,fill=BookPaper,line width=0.8pt,
    minimum width=24mm,minimum height=9mm,align=center,inner sep=2mm
  },
  relation/.style={
    draw=BookInk,line width=0.9pt,
    -{Latex[length=2mm]},rounded corners=1.2mm
  },
  edge label/.style={font=\figurefont\footnotesize,text=BookInk,inner sep=1mm}
]
  \node[anchor=west,font=\figurefont\small\bfseries,text=BookInk]
    at (0.1,4.65) {A\quad 输入覆盖};
  \node[anchor=west,font=\figurefont\small\bfseries,text=BookInk]
    at (8.1,4.65) {B\quad 有向任务收益};
  \fill[FigureInput!10] (0.5,0.8) rectangle (4.1,3.1);
  \fill[FigureOutput!10] (2.8,1.4) rectangle (4.1,3.1);
  \fill[BookAmber!12] (4.1,1.4) rectangle (6.8,3.6);
  \fill[BookAmber!12] (2.8,3.1) rectangle (4.1,3.6);
  \draw[BookBlue,line width=0.9pt] (0.5,0.8) rectangle (4.1,3.1);
  \draw[BookTeal,dashed,line width=0.9pt] (2.8,1.4) rectangle (6.8,3.6);
  \node[text=BookBlue] at (1.65,2.1) {来源};
  \node[text=BookTeal] at (5.55,3.2) {目标};
  \node[align=center,font=\figurefont\footnotesize,text=BookInk]
    at (3.45,2.2) {共同\\覆盖};
  \node[align=center,font=\figurefont\footnotesize,text=BookAmber]
    at (5.5,2.1) {目标中的\\来源缺口};
  \node[anchor=west,text=BookMuted,font=\figurefont\footnotesize]
    at (0.5,0.2) {重叠不代表标签机制相同};

  \node[task] (species) at (9.45,2.4) {物种识别};
  \node[task] (disease) at (14.7,2.4) {病害判断};
  \node[task] (localization) at (9.45,0.3) {花朵定位};
  \draw[relation] (species.north)
    -- (9.45,3.85) -- (14.7,3.85) -- (disease.north);
  \node[edge label,anchor=south] at (12.075,3.9) {$+2$};
  \draw[relation] (disease.west) -- (species.east);
  \node[edge label,anchor=south] at (12.075,2.45) {$-1$};
  \draw[relation] (localization.north) -- (species.south);
  \node[edge label,anchor=west] at (9.55,1.35) {$+4$};
  \node[align=left,anchor=west,text=BookMuted,font=\figurefont\footnotesize]
    at (11.45,0.3) {单位：百分点\\未画方向：未测};
\end{tikzpicture}

  \caption{来源覆盖与有向任务帮助关系回答不同问题。左侧是抽象输入区域的覆盖示意，面积不表示概率，重叠也不保证标签规则相容。右侧为教学构造，箭头从来源指向目标，标注相对该目标无迁移对照的准确率收益，单位为百分点；未画方向表示未测，不表示无效。}
  \label{fig:ket-reuse-coverage-relations}
\end{figure*}

这些箭头也不是共同训练关系。顺序迁移固定了来源学习结果，再让目标适应；共同训练时，两个任务会反过来改变共享表示、竞争参数容量并影响梯度。一个已经训练好的定位器有助于物种分类，并不能保证把定位损失加入物种模型的训练目标仍有同样收益。共同学习中的任务亲和性见第~\ref{sec:ket-sharing-affinity}~节，不能直接用这里的迁移箭头代替。

复用学习先验或适应规则时，关系还要提升到任务分布。沿用第~\ref{sec:ket-overview-learning-settings}~节的记号，以$\mathcal T_k$表示任务，$\Pi_{\mathrm{hist}}$表示历史任务的分布。每项任务规定输入输出接口及其语义、任务内分布$\mathcal P_k$和损失$\ell_k$，支持集与反馈的访问条件也须给定。一个从历史任务中学得的规则$U$把任务支持集$S_k$和既有状态变成任务预测器；简写$U(S_k)$时省略固定的既有状态。历史任务上的平均效果
\[
  \mathbb E_{\mathcal T_k\sim\Pi_{\mathrm{hist}}}
  [R_k(U(S_k))]
\]
只刻画该任务分布内的表现，其中$R_k$为任务$\mathcal T_k$的风险。若历史任务都是“每类数张照片的物种分类”，新任务改成病斑分割，支持集结构、输出维度和损失都变了，不能只因它们仍属植物图像就沿用快速适应的承诺。

应检查历史任务是否覆盖目标所需的变化，例如类别重组、光照变化和支持集规模，还是只在固定的类别规则下反复抽样。任务分布中的成功提供的是相应范围内的证据；目标所需变化完全未见时，应补做目标试训或保留不确定性。这些条件既约束来源选择，也约束第~\ref{sec:ket-transfer-task-splits}~节中的任务划分与学习策略准备。

\section{可复用知识的选择}
\label{sec:ket-reuse-selection}

适用依据把候选变得可解释，却未必足以排出顺序。候选很多时，把每个模型都完整适应一次可能比从头训练更贵；只看便宜分数，又可能漏掉慢热或互补的来源。本节按估计、选择与组合、检验与调整展开。候选既可以是整个模型，也可以是“某层表示加线性头”或“某来源中符合标注规范的样本子集”，评分对象必须与计划复用的对象一致。

\subsection{候选知识价值的估计}
\label{sec:ket-reuse-value-estimation}

估计应对应已经选定的收益口径。预测微调终点准确率、预测达到90\%所需标签、预测一次部署总成本，是三个不同目标。以下先讨论目标效果的估计，再把资源约束带回选择。固定模型上的代理评分省去目标网络试训，但仍可能使用目标标签和数值求解；目标试训则用更多计算观察适应动态，两者形成不同成本与证据范围的取舍。

\subsubsection{基于固定模型的价值估计}
\label{sec:ket-reuse-fixed-score}

当可以查询来源分类器的完整概率向量时，最直接的疑问是：它对目标照片的不同响应，能否稳定地对应目标类别？\term{对数期望经验预测}（log expected empirical prediction，LEEP）用目标标签建立这种经验对应，再用对应后的目标预测给来源评分。它不更新来源模型，也不要求源标签与目标标签一一同名，但要求来源有固定的离散概率接口、目标有带标签的分类样本\scite{ket-nguyen2020-leep}

设目标评分集为$S_{\mathrm{score}}=\{(\bm x_i,y_i)\}_{i=1}^{n}$，目标类别集合为$Y_T$，来源类别集合为$Y_S$。记$q_s(\bm x_i)$为来源给类别$s\in Y_S$的概率，满足$q_s\geq0$且$\sum_s q_s=1$。将每个目标样本的来源概率按其目标标签分组相加，得到
\begin{equation}
  \begin{aligned}
    \hat p(y,s)&=\frac1n\sum_{i:y_i=y}q_s(\bm x_i),\\
    \hat p(s)&=\sum_{y\in Y_T}\hat p(y,s),\qquad
    \hat p(y\mid s)=\frac{\hat p(y,s)}{\hat p(s)}.
  \end{aligned}
  \label{eq:ket-reuse-leep-joint}
\end{equation}
每个样本贡献总质量$1/n$，只是把这份质量软分给多个来源类别，因此$\hat p(y,s)$是归一化的经验联合分布。它描述来源输出与目标标签的对应，不是来源训练数据中的真实联合分布。若$\hat p(s)=0$，则评分集中所有样本的$q_s$均为零，该列对评分没有贡献，可跳过；不应执行除零。

对一张照片，先按来源概率考虑不同来源类别，再通过经验条件分布转成目标概率，即得到\term{期望经验预测器}（expected empirical predictor，EEP）：
\begin{equation}
  \begin{aligned}
    p_{\mathrm{EEP}}(y\mid\bm x)
    &=\sum_{s\in Y_S}\hat p(y\mid s)q_s(\bm x),\\
    s_{\mathrm{LEEP}}
    &=\frac1n\sum_{i=1}^{n}
      \log p_{\mathrm{EEP}}(y_i\mid\bm x_i).
  \end{aligned}
  \label{eq:ket-reuse-leep-score}
\end{equation}
这里和下文的$\log$均为自然对数。LEEP是不大于零的平均对数概率，越大表示这个经验预测器在评分集上越能给正确标签分配概率。零是完全预测的边界，不表示所有有限模型都能达到。形成条件分布与计算分数使用同一批标签，故它是代理估计，不是独立测试成绩。

表~\ref{tab:ket-reuse-score-data}~给出四张带病害标签的目标照片。候选A、B各有两个来源类别，另提供一维冻结特征。A的来源预测头可由特征$z_A$经过$q_1=\operatorname{sigmoid}(-\log(9)z_A)$得到；B的来源预测头恒输出$(0.5,0.5)$。两者都满足$q_2=1-q_1$，所以表中概率与特征来自相容的固定映射。

\begin{table}[htbp]
  \centering
  \small
  \begin{tabularx}{\linewidth}{rrXXrr}
    \toprule
    $i$ & $y_i$ & A的$(q_1,q_2)$ & B的$(q_1,q_2)$ & $z_{A,i}$ & $z_{B,i}$ \\
    \midrule
    1 & 0 & $(0.9,0.1)$ & $(0.5,0.5)$ & $-1$ & $-1$ \\
    2 & 0 & $(0.9,0.1)$ & $(0.5,0.5)$ & $-1$ & $1$ \\
    3 & 1 & $(0.1,0.9)$ & $(0.5,0.5)$ & $1$ & $-1$ \\
    4 & 1 & $(0.1,0.9)$ & $(0.5,0.5)$ & $1$ & $1$ \\
    \bottomrule
  \end{tabularx}
  \caption{LEEP与LogME共用的教学评分数据。$y=0$表示健康、$y=1$表示病害；$z_{A,i}$和$z_{B,i}$是标量特征，非真实图像测量。}
  \label{tab:ket-reuse-score-data}
\end{table}

按行对应目标类别0、1，按列对应来源类别1、2，A的联合分布和条件分布分别为
\[
  \bigl[\hat p_A(y,s)\bigr]=
  \begin{pmatrix}0.45&0.05\\0.05&0.45\end{pmatrix},
  \qquad
  \bigl[\hat p_A(y\mid s)\bigr]=
  \begin{pmatrix}0.9&0.1\\0.1&0.9\end{pmatrix}.
\]
每个来源列的边缘概率都是0.5。对第一张照片，EEP给正确标签0的概率为$0.9\times0.9+0.1\times0.1=0.82$；另外三张的正确标签概率也都是0.82，因此$s_{\mathrm{LEEP},A}=\log0.82\approx-0.198451$。B的联合表四格均为0.25，条件表四格均为0.5，得分为$\log0.5\approx-0.693147$。按LEEP应把A排在B之前。

这个计算还给出几个边界。若来源对所有照片输出相同概率，EEP只能恢复目标类别频率，LEEP便等于负的经验标签熵；类别越不平衡，这个数也可能越接近零，却不代表稀有病害已经能识别。因此应在同一个目标集上比较来源，不能把不同任务的原始分数当成统一能力尺度。只返回一个硬类别的接口丢失了软概率信息；把它替换成独热向量可以另作评分，但不等价于原分类器的完整概率版LEEP。

LEEP原文还给出一个有条件的关系：若冻结表示上的候选预测头集合包含EEP，那么该集合最优的训练平均对数似然至少等于LEEP，因为EEP本身就是可选解。这不是对独立测试准确率的下界，也不能无条件套到不包含EEP的线性头集合。实际头能否优化到足够好的解、是否泛化，以及解冻后表示怎样变化，都需要另外检查。

有些候选只提供编码器，没有来源分类头；目标也可能是花朵位置等连续量。这时可以问：固定特征能否在一个受约束的预测模型中解释目标标签？\term{最大证据的对数}（logarithm of maximum evidence，LogME）用贝叶斯线性模型回答这个问题，给所有可能的线性权重按先验加权后计算标签证据，再优化少量超参数。概率边缘化的通用含义见第~\ref{sec:pgm-parameter-learning-bayesian}~节，这里展开评分所需的具体计算\scite{ket-you2021-logme}

从某个来源提取特征并按行组成$\bm Z\in\mathbb R^{n\times d}$。先考虑一维实值标签$\bm y\in\mathbb R^n$，假设
\begin{equation}
  \begin{aligned}
    \bm w&\sim\mathcal N(\bm0,\alpha^{-1}\bm I_d),\\
    \bm y\mid\bm Z,\bm w
    &\sim\mathcal N(\bm Z\bm w,\beta^{-1}\bm I_n),
    \qquad \alpha,\beta>0.
  \end{aligned}
  \label{eq:ket-reuse-logme-model}
\end{equation}
其中$\alpha$控制权重先验的精度，$\beta$控制独立观测噪声的精度。先验偏好较小权重，但不是把权重固定为零；噪声项允许线性模型无法完全解释标签。这里使用无显式截距的原式形式，若另外加入常数特征、中心化或标准化，都要把处理规则固定并记录，因为它们会改变评分模型。

只最大化$p(\bm y\mid\bm Z,\bm w)$可能用很大权重拟合偶然噪声。模型证据则把权重积分掉：
\begin{equation}
  p(\bm y\mid\bm Z,\alpha,\beta)
  =\int p(\bm y\mid\bm Z,\bm w,\beta)
       p(\bm w\mid\alpha)\,\mathrm d\bm w.
  \label{eq:ket-reuse-logme-evidence}
\end{equation}
令$\bm A=\alpha\bm I_d+\beta\bm Z^\top\bm Z$，$\bm m=\beta\bm A^{-1}\bm Z^\top\bm y$。完成平方可以看到
\[
  \begin{aligned}
    &\alpha\lVert\bm w\rVert_2^2+
      \beta\lVert\bm y-\bm Z\bm w\rVert_2^2\\
    &\quad=(\bm w-\bm m)^\top\bm A(\bm w-\bm m)
       +\alpha\lVert\bm m\rVert_2^2
       +\beta\lVert\bm y-\bm Z\bm m\rVert_2^2.
  \end{aligned}
\]
对配方后的高斯指数因子$\exp\{-\tfrac12(\bm w-\bm m)^\top\bm A(\bm w-\bm m)\}$关于$\bm w$积分，得到$(2\pi)^{d/2}|\bm A|^{-1/2}$；代回归一化常数后，得到原论文的对数证据：
\begin{equation}
  \begin{aligned}
    L(\alpha,\beta)
    &=\frac n2\log\beta+\frac d2\log\alpha
       -\frac n2\log(2\pi)-\frac12\log|\bm A|\\
    &\quad-\frac\beta2\lVert\bm y-\bm Z\bm m\rVert_2^2
       -\frac\alpha2\lVert\bm m\rVert_2^2.
  \end{aligned}
  \label{eq:ket-reuse-logme-log-evidence}
\end{equation}
残差项衡量解释标签的程度，权重项和行列式项共同反映先验与有效模型复杂度。模型证据因此不同于训练误差，但这种复杂度处理仍依赖线性、高斯和各向同性先验等假设，不能把它说成任何数据和模型下都不会过拟合的保证。

评分还需选择$\alpha,\beta$。设$\lambda_j$为$\bm Z^\top\bm Z$的非负特征值，在第$t$次迭代用当前超参数计算后验均值$\bm m^{(t)}$及有效维数$\gamma^{(t)}$，再更新
\begin{equation}
  \begin{aligned}
    \gamma^{(t)}&=\sum_{j=1}^{d}
       \frac{\beta^{(t)}\lambda_j}{\alpha^{(t)}+\beta^{(t)}\lambda_j},\\
    \alpha^{(t+1)}&=
       \frac{\gamma^{(t)}}{\lVert\bm m^{(t)}\rVert_2^2},\\
    \beta^{(t+1)}&=
       \frac{n-\gamma^{(t)}}{\lVert\bm y-\bm Z\bm m^{(t)}\rVert_2^2}.
  \end{aligned}
  \label{eq:ket-reuse-logme-update}
\end{equation}
这里的有效维数按数据相对于先验的约束强度累计各特征方向：$\beta^{(t)}\lambda_j$远大于$\alpha^{(t)}$时该方向接近计为一维，远小时接近零维。上述三式给出证据最大化的固定点更新，不是对来源网络执行梯度训练。原算法从$\alpha=\beta=1$开始，利用特征矩阵分解重复计算$\bm m$和$\gamma$。实际实现应限定最大迭代数，并按超参数或证据的相对变化停止；零残差、零权重范数及边界最优需要专门处理，不能机械执行除零，也不能把有限次迭代一律称为已取得全局最大值。

对$c$类分类，原方法将标签转成$n\times c$的独热矩阵$\bm Y$，对每一列$\bm Y_{:,j}$分别使用上述高斯回归模型、分别求解超参数，再平均每样本的最大对数证据：
\begin{equation}
  s_{\mathrm{LogME}}
  =\frac1c\sum_{j=1}^{c}\frac1n
    \sup_{\alpha_j,\beta_j>0}
    L_j(\alpha_j,\beta_j).
  \label{eq:ket-reuse-logme-score}
\end{equation}
多输出回归同样逐列处理。分类时，这是把独热标签作为回归目标的代理模型，不是精确计算softmax分类器的边际似然。分数越大越受该模型支持；由于它是连续密度的对数，可能为正，不能沿用LEEP不大于零的范围。回归标签的单位、特征预处理及目标集合应保持一致。这里写上确界，是因为最优证据可能只在$\alpha$或$\beta$趋于边界时逼近而没有有限最大点；数值实现应报告搜索约束及边界解。

用表~\ref{tab:ket-reuse-score-data}~继续复算。对病害标签列$\bm y=(0,0,1,1)^\top$，A的特征列$\bm z_A=(-1,-1,1,1)^\top$满足$\bm z_A^\top\bm z_A=4$、$\bm z_A^\top\bm y=2$、$\bm y^\top\bm y=2$。为看清最优超参数，可将证据等价地写成零均值高斯密度，协方差为
\[
  \bm C=\beta^{-1}\bm I_4+\alpha^{-1}\bm z_A\bm z_A^\top.
\]
沿单位向量$\bm u=\bm z_A/2$，标签的平方投影为$(\bm u^\top\bm y)^2=1$，其余三个正交方向合计平方长度为1。记$r=\beta^{-1}$、$v=r+4/\alpha$，则
\[
  L_A=-\frac12\left[
  4\log(2\pi)+\log v+3\log r+\frac1v+\frac1r\right].
\]
分别求导得到$v=1$、$r=1/3$，符合$v>r$，所以$\beta=3$、$\alpha=6$。此时$\bm m$是标量$1/3$，残差平方和为$10/9$，代入式~\eqref{eq:ket-reuse-logme-log-evidence}~可得
\[
  \frac{L_A}{4}
  =-\frac{4\log(2\pi)-3\log3+4}{8}
  \approx-1.006959.
\]
健康标签列$(1,1,0,0)^\top$只改变投影符号，证据相同，因此这也是A的两类平均LogME。固定点迭代的前两次超参数为$(5,3.076923)$和$(5.625,3.035552)$，与解析最优点$(6,3)$的差距逐步缩小。

B的特征列$\bm z_B=(-1,1,-1,1)^\top$与两列独热标签都正交。沿$\bm z_B$增加方差不能解释标签，只会增加体积惩罚，证据最优趋于$\alpha\to\infty$，即放弃这个特征；其余最优为$\beta=2$。故
\[
  s_{\mathrm{LogME},B}
  =-\frac12\bigl(\log\pi+1\bigr)
  \approx-1.072365.
\]
这里没有有限的最优$\alpha$，应报告边界解；设置一个很大的有限上限只是数值近似。A仍排在B之前，但两种评分的绝对数值和间隔没有共同尺度，不能把$-0.198451$与$-1.006959$相减评价方法优劣。

两种评分能看到的候选范围也不同。若B只开放预测接口，其LEEP可算，LogME不可算；若某候选只开放自监督编码器、没有固定来源类别概率，其LogME可算，原始LEEP则缺少输入条件。对同一目标评分集，LEEP的额外计算主要是累计来源概率和目标类别的对应；LogME需要保存或分批处理特征并做矩阵分解。直接构造$\bm Z^\top\bm Z$及分解的成本约为$O(nd^2+d^3)$，后续各标签列可复用分解，避免反复求逆；$d$很大时还要考虑秩和存储，不能只说“不微调所以没有成本”。

论文通过候选分数与实际迁移效果之间的相关性检验评分，但排序相关和收益本身是两份证据。例如三个候选的代理排名为A、B、C，充分适应后的排名为B、A、C，无并列时三对候选中两对一致、一对相反，Kendall秩相关为$(2-1)/3=1/3$；代理第一名仍选错了。即使排名完全正确，所有候选也可能都不如无迁移对照。因此应分别报告候选集合、评分标签量、排序质量和相对对照的真实收益，不能用一个相关系数替代最终判断。

\subsubsection{基于目标试训的价值估计}
\label{sec:ket-reuse-target-trials}

固定特征上的代理模型看不到解冻后表示如何变化，也看不到实际优化器是否适合某个来源。\term{目标试训}（target-side trial training）用受控的小预算运行目标学习过程，观察候选在计划使用方式下的表现。它不是把完整训练随意缩短，而是明确试训对象、预算和可用于选择的证据。

冻结表示只训练轻量头，主要检验信息能否被该头读出；只更新少量层或少量参数，还能观察局部调整能否修正目标失配；完整适应才能覆盖全部参数协同变化和较长优化动态，但成本最高。三者观察的对象不同。一个在线性探测中低分的来源，可能解冻后有潜力；一个冻结时高分的来源，也可能在小样本全量微调中丢掉有效结构。试训结论应连同更新范围记录。

可将目标开发数据分为用于拟合的$S_{\mathrm{fit}}$和用于候选决策的$S_{\mathrm{select}}$。候选使用相同的样本子集、类别组成、预处理、批量顺序、训练步数和验证频率；为初始化、优化器、学习率及正则指定相同的规则和调参机会。固定同一个学习率有利于控制变量，却可能偏向某类模型，所以若最终要比较各自适配后的能力，应在同样大小的预设学习率网格内选择，而不是对一个来源精调、对另一个只跑默认值。优化步数控制不住模型大小差异时，再增加设备时间或运算量上限。

可用表~\ref{tab:ket-reuse-score-data}~的四个特征样本具体计算一次冻结头试训。将它们用于拟合，另取四张不参与拟合的目标照片作为选择集；教学上设这四张新照片恰好具有相同的特征与标签组合，并非把训练照片再次充作验证集。把二元标签编码为$t_i=2y_i-1$，使用无截距的标量头$wz_i$，全批量最小化
\begin{equation}
  H(w)=\frac{1}{2n}\sum_{i=1}^{n}(wz_i-t_i)^2,
  \qquad
  w^{(t+1)}=w^{(t)}-\eta H'(w^{(t)}).
  \label{eq:ket-reuse-trial-update}
\end{equation}
这里$t_i$是样本的标签编码，上标$(t)$是优化步，两者作用不同。固定$w^{(0)}=0$、$\eta=1/2$，不加正则，恰好训练两步。对A，$n^{-1}\sum_i z_i^2=1$且$n^{-1}\sum_i z_it_i=1$，故$H'_A(w)=w-1$，依次得到$w_A^{(1)}=1/2$、$w_A^{(2)}=3/4$。选择集上的$H_A$依次为$1/8$、$1/32$；对B，$n^{-1}\sum_i z_it_i=0$，因此$H'_B(w)=w$，两步权重均为0，选择损失均为$1/2$。按相同两步预算，应选择A的冻结表示。

使用时按$\mathbb I(wz\geq0)$给出类别，并规定平局归为1。A在这组选择照片上两步都达到100\%准确率，B为50\%，说明试训损失继续下降与分类准确率继续上升也不是同一件事。这次试训没有更新特征，只证明指定预测头和预算下A更可用；若允许B的编码器更新，原先的计算不再能预测其终点表现。

对较长试训，在第$t$步记录$\hat R_{\mathrm{select}}(f_j^{(t)})$及重复运行的波动，而不只保存最好的一个数。另考虑一个允许更新来源参数的教学构造，两个候选在100张选择照片上的错误数如下：A在第2、8、16步分别错20、18、18张，B分别错22、16、14张。若只跑2步就淘汰B，将选择80\%而非78\%的早期准确率，却错过第16步86\%而非82\%的方案。这和第~\ref{sec:ket-reuse-prediction}~节的起点与终点区别相同，只是现在它直接影响搜索。

多阶段预算可以减少完整试训次数。下面给出一个明确的教学流程：8个候选各花1个计算单位做固定模型评分，保留4个；每个保留候选试训2步，保留2个；这2个候选都从检查点继续到8步，再继续到16步后比较。假设一次试训步的计算恰为1单位，且候选间相同，则总成本为
\[
  8\times1+4\times2+2\times(8-2)+2\times(16-8)
  =44.
\]
所有8个候选从头完整试训16步的成本为128单位。这两个数都包括其各自采用的筛选计算；被淘汰候选已消耗的步骤没有从账本中抹去。若不保留检查点而重新训练，应改计全部重训步骤，不能仍使用增量差。

在这个流程中，早期筛选仍可能漏掉好来源。保留数量越少，成本越低，误排风险通常越需要关注；应为不同接口、不同表示类型或不确定性较大的候选保留少量探索预算。对于预先固定的同一训练进度，若候选$j$的风险区间$[\underline r_j,\overline r_j]$满足$\underline r_j>\min_h\overline r_h+\delta$，可以把它列为淘汰对象。但若算法根据已见结果自适应地决定下一次查看、重复检验同一候选或提前停止，逐阶段普通置信区间不再自动控制整个淘汰过程的误差；须使用覆盖全部允许查看时刻的序贯有效区间，或预先固定阶段并对多次比较作校正。即使区间在当前步有效，也不能保证终点风险满足同样关系。没有可靠的学习曲线外推条件时，这只是预算策略，不是无误淘汰定理。

已经反复用来调学习率、看曲线和选来源的$S_{\mathrm{select}}$承担的是开发工作。无论称它为验证集还是测试集，它都不再是独立的最终证据。确认候选后，应锁定来源、更新范围、训练预算和停止规则，再在未参与这些决策的数据上评价。第~\ref{chap:ket-evaluation}~章进一步讨论重复使用选择数据造成的偏差。

\subsection{候选知识的选择与组合}
\label{sec:ket-reuse-choice}

估计分数回答某个候选可能多有价值，选择还要回答能否实际使用、是否值得付出代价，以及需要多少来源。先把必须满足的接口、授权、存储与延迟写成约束，再在可行候选中比较收益。这样得到的不是脱离场景的模型排行榜，而是一个目标和预算下可执行的复用方案。

\subsubsection{单来源知识的选择}
\label{sec:ket-reuse-single-source}

令$J$为来源候选索引集合，并增加$j=0$表示无迁移方案。记$J_0=J\cup\{0\}$，$\hat r_j$为候选$j$在既定训练终点的预计目标风险，$c_j$为选择后仍需支付的适应成本，$m_j$为部署存储，$a_j$表示所需接口和授权是否可用；无迁移方案的这些量也按同一核算边界实测，不把它的训练成本设为零。在预算$B_c,B_m$下，一个基本选择形式是
\begin{equation}
  \begin{aligned}
    j^\star&\in\argmin_{j\in J_0}\hat r_j\\
    \text{满足}\quad
    &c_j\leq B_c,\quad m_j\leq B_m,\quad a_j=1.
  \end{aligned}
  \label{eq:ket-reuse-single-choice}
\end{equation}
这里的$\hat r_j$应来自相同目标损失、数据协议与训练终点下的风险估计，例如预算匹配的目标试训及其验证结果。LEEP、LogME等代理分数可用于安排试训次序，但其排名不能直接换算成准确率或填入$\hat r_j$；若要据代理预测终点风险，还需用独立的迁移记录建立并检验这种对应。

若所有可行来源的预计风险都高于无迁移方案，式~\eqref{eq:ket-reuse-single-choice}~会选择$j=0$，不会为了“必须迁移”而返回一个有害来源；若连无迁移方案也不满足硬约束，则应报告当前没有可执行方案，而不是自动放宽约束。若目标优先节省标签，就把固定效果下的标签预算作为目标，并同样保留无迁移候选；若存在多个不能互换的目标，可保留效果与成本都不被另一候选同时压倒的方案，再由应用要求取舍。

假设候选A在LEEP上最高，却只有远程预测接口，而计划是取得参数后在离线设备上微调和部署，那么A的评分并不能把它变成可执行方案。候选B预计准确率90\%、存储400 MB，候选C预计89.8\%、存储80 MB；在100 MB部署上限下，B不可行，C才进入比较。这个教学对照同时表明，代理分数计算所需接口可能弱于最终部署所需接口，筛选之前就应核对两者。

若B、C都可部署，预计风险的差又小于估计不确定性，不应把小数点后的顺序当作可靠结论。可以增加配对试训、验证更关键的子群，或暂时保留两个候选；在硬预算下必须立即作决定时，应写明证据不足，并按较低成本或较小风险暴露等预设次级规则处理。不能事后从多个次级标准中挑选恰好支持偏好模型的一个。

单来源不等于照搬全部内容。来源确定后，还可排除标注标准冲突的样本，选取保留目标纹理的中间层，或只使用定位模块而不保留其来源分类头。每一次选择都应落实到可识别的对象，例如模型版本、层名、样本纳入条件和模块输入输出。删除样本要同步检查目标覆盖是否变差；只保留较早层也可能增加目标端重建高层结构的负担。怎样冻结、重训或增加任务参数由第~\ref{chap:ket-transfer-adaptation}~章展开，本节需要留下的是对象与证据对应的选择记录。

\subsubsection{多来源知识的选择}
\label{sec:ket-reuse-multiple-sources}

当一个来源不能覆盖目标要求时，可以选择多个来源。但两个模型都很准确，不代表它们放在一起更有用：它们可能在同一批照片上出错，也可能提供互补的信息，还可能因标签规范或预测校准冲突而互相破坏。选择要关注加入一个来源以后增加了什么，而不是把各自成绩直接相加。

Taskonomy用分别训练的任务编码器建立来源到目标的迁移关系，再把这些关系用于监督预算约束下的来源与路径选择。其迁移函数在冻结来源表示上用较少目标数据训练低容量读取器；高阶迁移同时读取多个来源表示，用于测量组合可能带来的互补。论文先在各目标内部归一化迁移关系，再用二元整数规划选择来源和进入目标的迁移路径。这是特定任务字典、编码器、读取器和监督预算下的有向关系，不是所有任务天然固定的距离，也不是共同训练时的梯度相容性\scite{ket-zamir2018-taskonomy}

一个与这种思路相符的选择模型，可以把一条候选路径$e$表示为$(J_e,k_e)$：使用来源集合$J_e$服务目标$k_e$。令$u_j\in\{0,1\}$表示是否选中来源$j$，$v_e\in\{0,1\}$表示是否选中路径$e$，$\hat g_e$为该路径在统一目标效用口径下的预计收益，$w_k$为目标重要性，$c_j$和$d_e$分别为来源与路径成本。则可写为
\begin{equation}
  \begin{aligned}
    \max_{\bm u,\bm v}\quad
      &\sum_e w_{k_e}\hat g_ev_e\\
    \text{满足}\quad
      &v_e\leq u_j \quad(j\in J_e),\\
      &\sum_{e:k_e=k}v_e=1 \quad(\text{每个目标 }k),\\
      &\sum_jc_ju_j+\sum_ed_ev_e\leq B.
  \end{aligned}
  \label{eq:ket-reuse-path-choice}
\end{equation}
每个目标的路径集合中可加入无来源基线，避免模型强迫所有目标迁移。该式是本章为解释选择而写的成本扩展形式，不逐项复现Taskonomy原文约束；组合收益$\hat g_e$必须来自组合试训或有依据的估计，而不是单来源得分求和。整数规划求得的最优解也只对所列候选、估计值和约束成立，不能自动成为真实效果的全局最优解。

表~\ref{tab:ket-reuse-complementarity}~给出三个来源、两个目标的完整教学构造。目标$\mathcal T_1$为物种识别，$\mathcal T_2$为病害判断，各有100个等权输入组，每组放入一个验证对象；组号是可观察的输入属性，不是验证样本的行号。各自无迁移对照恰好在组1至20上出错。

为让组合可复算，规定来源提供的是带可用标志的修正模块：在指定输入组上给出修正，其他组拒绝介入。可用标志只依赖组号，不查看真实标签；本例人为安排所有被触发的修正都正确，重叠修正也一致。因而组合规则是在任一模块提供修正时采用它，否则保留基线，组合修正范围恰为集合并集。这是构造的机制，不是从单模型准确率推断出的真实集成效果。

具体地，A在两个目标上都修正组$\{1,\ldots,6\}$；B在$\mathcal T_1$上修正$\{1,\ldots,8\}$，在$\mathcal T_2$上只修正$\{9\}$；C在$\mathcal T_1$上只修正$\{9\}$，在$\mathcal T_2$上修正$\{1,\ldots,8\}$。每个来源的加载与适应成本为1个教学预算单位，组合额外成本在此假设为零，总预算为2；两个目标等权。所有结果都可由上述集合直接计数，来源名称不指代真实模型。

\begin{table*}[htbp]
  \centering
  \small
  \begin{tabularx}{\linewidth}{lrrrrX}
    \toprule
    来源集合 & $\mathcal T_1$收益 & $\mathcal T_2$收益 & 平均收益 & 总成本 & 组合范围与预算状态 \\
    \midrule
    $\varnothing$ & 0 & 0 & 0 & 0 & 无迁移对照，各错20张 \\
    $\{A\}$ & 6 & 6 & 6 & 1 & 两个目标都修正组1至6 \\
    $\{B\}$ & 8 & 1 & 4.5 & 1 & 偏向物种目标 \\
    $\{C\}$ & 1 & 8 & 4.5 & 1 & 偏向病害目标 \\
    $\{A,B\}$ & 8 & 7 & 7.5 & 2 & 物种修正大部重叠 \\
    $\{A,C\}$ & 7 & 8 & 7.5 & 2 & 病害修正大部重叠 \\
    $\{B,C\}$ & 9 & 9 & 9 & 2 & 两个目标分别补上组9 \\
    $\{A,B,C\}$ & 9 & 9 & 9 & 3 & 超出预算，且A不增加覆盖 \\
    \bottomrule
  \end{tabularx}
  \caption{三个来源、两个目标的互补反例，全部为教学构造。收益单位为准确率百分点，成本为预设预算单位；组合成绩由正文给定的可用标志和修正规则计算，不是论文测量。在实际选择中，表内每个组合的对应值均须单独验证。}
  \label{tab:ket-reuse-complementarity}
\end{table*}

只按单项平均收益排序，会先选A，再选B或C，最终平均收益7.5个百分点。枚举预算内的全部集合，却得到$\{B,C\}$的9个百分点，两个目标的错误数均从20降为11。以$\{A,B\}$为例，$\mathcal T_1$的修正并集只有$\{1,\ldots,8\}$，不能把A的6与B的8加成14；$\mathcal T_2$的并集是$\{1,\ldots,6,9\}$，得到7。单来源价值较低的B、C恰好补齐彼此的弱项。

即使用真实边际收益逐步贪心，第一步仍会选A，第二步最多得到7.5，因此这个例子也不支持“每一步挑当前最好就必得最优组合”。这里的并集计数具有覆盖函数的递减边际性质，但一般模型组合未必如此：加入来源可能导致冲突、挤占容量或改变路由，使原有正确预测变错，甚至破坏单调性。没有明确的集合函数性质和成本条件，就不能套用贪心近似保证。对少量候选可直接枚举重要组合，对大量候选则应保留覆盖不同目标区域的来源，再为候选组合分配验证预算。

本节决定的是选哪些来源及哪些组合值得检验。真实模块怎样共同计算、预测如何融合、参数能否合并、能否蒸馏成单模型，由第~\ref{chap:ket-integration}~章展开。若改成让来源任务共同学习，则选择对象与相互影响又发生变化，应转到第~\ref{chap:ket-sharing}~章分析，不能直接沿用表中的固定组合值。

\subsection{选择结果的检验与调整}
\label{sec:ket-reuse-validation}

候选被选中以后，原先的适用依据和估计分数仍是可被推翻的判断。应按预定口径检查目标效果、阈值到达量、实际成本及必须保留的能力，尤其查看此前指出的缺口区域。整体准确率提高，而新增物种全部被归入旧类，可能没有达到真实的目标要求；最终训练较快，但来源搜索耗时更长，也可能没有实现资源收益。

发现有害结果时，应把原因分开检验。来源缺少目标相关信息时，冻结探测、充分训练的头和解冻试验都可能暴露同一覆盖缺口；代理误排时，分数高的来源在匹配预算的目标试训中落后；适应方法不合适时，同一个来源改变学习率、冻结范围或标签映射后可能恢复效果。这些迹象帮助提出假设，却不互相构成证明，应通过只改变相关因素的对照确认，而不是把所有失败归为“来源模型差”。

调整可以是重新选来源、只保留有依据的层或样本、把冲突能力放入独立分支，或停止迁移并回到无迁移对照。后两种选择并非研究失败：它们让系统在证据不足或成本不合算时避免继续扩大风险。每次调整都要重新登记目标标签、无标签输入、试训轮次、模型版本和累计计算；重新排列同一批分数本身不会带来新证据。

图~\ref{fig:ket-reuse-selection-loop}~把选择与验证放在同一个开发闭环中。候选集合和访问条件限制可估计对象，价值证据可以来自固定模型评分，也可以直接来自预算试训；评分不是必经步骤。单来源与多来源在选择节点分开，随后对所选方案继续试训与复核。预设停止条件应包括达到效果与资源要求、剩余预算耗尽、没有可行候选或继续试训已无法解决当前不确定性。不能把“直到验证分数满意”为止当成无代价的停止规则。

\begin{figure*}[htbp]
  \centering
  % Source file only. Development control flow is solid; evaluation data flow is dashed.
% All centerline connectors use shared x/y coordinates and explicit node anchors.
% Final evaluation has no return edge to development.
\begin{tikzpicture}[
  x=1cm,y=1cm,font=\figurefont\small,
  box/.style={
    rectangle,draw=BookInk,fill=BookPaper,line width=0.8pt,
    text width=42mm,minimum height=9mm,
    align=center,inner xsep=2mm,inner ysep=1.5mm
  },
  branch/.style={box,text width=28mm,fill=FigureModel!10},
  flow/.style={
    draw=BookInk,line width=0.9pt,rounded corners=1.2mm,
    -{Latex[length=2mm]}
  },
  edge label/.style={
    font=\figurefont\footnotesize,text=BookInk,inner sep=1mm
  }
]
  \node[box,fill=FigureInput!10,text width=50mm] (input)
    at (6,0) {候选集合与访问条件};
  \node[box] (score) at (6,-1.45) {候选价值估计\\固定评分或预算试训};
  \node[box] (choose) at (6,-2.9) {预算约束下选择};
  \node[branch] (single) at (3.2,-4.4) {单来源\\模型、层或子集};
  \node[branch] (multiple) at (8.8,-4.4) {多来源\\组合与路径};
  \node[box,fill=FigureModel!10] (trial)
    at (6,-5.95) {目标试训};
  \node[box] (review) at (6,-7.4) {判断：效果、成本\\与保留要求是否满足};
  \node[box,fill=FigureOutput!12] (lock)
    at (6,-9.3) {接受并锁定方案};
  \node[box,text width=33mm,fill=BookAmber!10] (stop)
    at (12.65,-5.95) {终止或回到对照};
  \node[box,text width=33mm,dashed,fill=white] (testdata)
    at (12.65,-7.7) {独立最终评价数据\\开发期间封存};
  \node[box,text width=33mm,fill=FigureOutput!12] (evaluate)
    at (12.65,-9.3) {最终评价与报告};

  \draw[flow] (input.south) -- (score.north);
  \draw[flow] (score.south) -- (choose.north);
  \draw[BookInk,line width=0.9pt] (choose.south) -- (6,-3.6);
  \draw[flow] (6,-3.6) -- (3.2,-3.6) -- (single.north);
  \draw[flow] (6,-3.6) -- (8.8,-3.6) -- (multiple.north);
  \draw[BookInk,line width=0.9pt,rounded corners=1.2mm]
    (single.south) -- (3.2,-5.25) -- (6,-5.25);
  \draw[BookInk,line width=0.9pt,rounded corners=1.2mm]
    (multiple.south) -- (8.8,-5.25) -- (6,-5.25);
  \draw[flow] (6,-5.25) -- (trial.north);
  \draw[flow] (trial.south) -- (review.north);
  \draw[flow] (review.south) -- (lock.north);
  \node[edge label,anchor=west] at (6.1,-8.55) {满足};
  \draw[flow] (lock.east) -- (evaluate.west);
  \draw[flow,dashed] (testdata.south) -- (evaluate.north);

  \draw[flow] (choose.east) -- (12.65,-2.9) -- (stop.north);
  \node[edge label,anchor=west] at (12.8,-4.35) {无可行候选};
  \draw[flow] (review.east) -- (9.55,-7.4)
    -- (9.55,-5.95) -- (stop.west);
  \node[edge label,anchor=west,align=left]
    at (9.65,-6.7) {不满足\\且不可继续};
  \draw[flow] (review.west) -- (0.4,-7.4)
    -- (0.4,-1.45) -- (score.west);
  \node[edge label,anchor=west,align=left]
    at (0.55,-6.45) {不满足且可继续\\调整候选或范围\\登记新增预算};
  \node[text=BookMuted,font=\figurefont\footnotesize]
    at (12.65,-10.15) {最终数据不返回选择环节};
\end{tikzpicture}

  \caption{候选估计、受约束选择与目标试训构成开发闭环。候选估计按接口与预算选择固定评分、试训或二者结合；两个来源分支共享后续试训与复核。只有复核未通过且预算仍允许时才调整。最终评价数据处于循环之外，在决策锁定后使用；若据其结果再次调整，它就已成为开发数据。}
  \label{fig:ket-reuse-selection-loop}
\end{figure*}

接受候选时，应锁定来源标识、复用范围、预处理、目标训练与停止规则，再进行独立最终评价。预算耗尽而证据仍不确定时，应报告“不足以确认收益”，而不是把未检出负迁移当成证明安全。若最终评价暴露新的问题，可以开启下一轮研究，但原评价数据已被用于决策，下一轮独立证据需要重新安排。选择、适应和评价的职责由此分清，第~\ref{chap:ket-evaluation}~章将给出完整的实验与统计协议。

\FloatBarrier
\section{本章小结}
\label{sec:ket-reuse-summary}

面对新的花卉目标任务，不能仅凭来源规模、来源准确率或照片相似度决定复用。应先明确要改善预测、节省目标样本，还是减少整个使用期限内的成本，再核对输入输出、表示和任务结构的适用条件。早期学习更快与最终风险更低可以有相反答案，一次目标适应便宜与包含来源准备的总成本更低也可以有相反答案。

LEEP用来源概率与目标标签的经验对应估计候选价值，LogME用冻结特征下的线性高斯模型证据估计价值；两者都能减少逐个完整适应的负担，但都有访问、标签和建模条件。目标试训补充适应动态的证据，仍可能受早期误排和选择数据反复使用的影响。可行性约束决定哪些高分候选真正能用，多来源的互补、重复与冲突则要求比较组合效果，不能只累加单项分数。

选择的结果应是一份带条件的决定：使用什么、为何适用、相对什么对照预期获得什么收益、花费多少，以及什么证据会触发调整或停止。负迁移由相对对照的损失确认，覆盖不足、分数误排和适应不当是需要进一步辨认的原因。这样的决定把后续更新、整合、迁移与共享建立在可检验的依据上，也保留了证据不够时不复用的选项。
